\documentclass{homework}

\usepackage{enumitem}
\usepackage{tcolorbox}

\usepackage{tikz}
\usetikzlibrary{positioning,calc}

\usepackage{listings}

\lstset{
basicstyle=\small\ttfamily,
columns=flexible,
breaklines=true
}

\usepackage{fancyvrb}
\fvset{listparameters=\setlength{\topsep}{0pt}\setlength{\partopsep}{0pt}}

\title{04}

\begin{document}
\maketitle

\textbf{Organisational details}

\begin{itemize}
  \item Include in your report approximately how long the homework took you.
  \item You must follow the submission instructions outlined on GitLab.
  \item If the GitLab pipeline does not succeed, I will not grade your assignment (you will receive 0).
  The same applies if the Markdown report is illegibly formatted or if you miss the final submission deadline.
  \item I will provide feedback on Discord and/or GitLab.
  Check both after receiving your grade.
  \begin{itemize}
    \item My initial feedback may be quite minimal as there are so many of you.
    However, you can request additional feedback.
    Specific questions are also welcome.
    \item Do not share the feedback you receive, as this may give others an unfair advantage.
  \end{itemize}
  \item Grace clause submission deadline: 19.04.2026 23:59 EEST.
  \begin{itemize}
    \item The grace clause applies only if you submit the full work (all tasks) by the soft deadline and the GitLab pipeline succeeds.
    \item You have one week from receiving the feedback to fix and resubmit your improved work.
    You may regain up to (but not necessarily) half of the lost points.
  \end{itemize}
  \item You may share hints with each other, but not answers or code.
  Ask questions in the server if they could benefit everyone.
  \item You may -- and I encourage you to -- ask me questions.
  \item You may use AI as assistance, but not to solve the tasks themselves.
  \begin{itemize}
    \item Acceptable: e.g. checking your understanding, asking how to loop in Python.
    \item Not acceptable: e.g. generating code logic, library use, or crypto operations.
    \item If you use AI in any form, including for rewording, state in your report where and what you used it for.
    \item I recommend asking in the server or contacting me instead of using AI.
  \end{itemize}
  \item If I catch you plagiarising, copying, or using AI for code logic, I will report you to the faculty and fail you for the course.
  \begin{itemize}
    \item If during the exam you are unable to explain material for which you received homework credit, I may ask you to explain your work.
    \item If you cannot convince me that the homework is your own work, I will report you to the faculty and fail you.
  \end{itemize}
\end{itemize}

\newpage

\begin{center}
  Theory tasks
\end{center}

\textbf{Additional instructions}

\begin{itemize}
  \item \textbf{Please do not use AI for the theoretical part.}
  Challenge your understanding and look for additional written sources if the lecture and provided materials are not enough.
  Critical thinking and clear communication are skills -- use them or lose them.
  \item You must reference all external materials you rely on, including AI tools, if you choose to use them anyway.
\end{itemize}

\begin{task}{GPG}
  We briefly covered the web-of-trust approach and the existence of PGP, however, we did not go in depth on what PGP is.

  PGP stands for \emph{Pretty Good Privacy} and is a cryptographic toolset which can be used for encryption and signing of data.
  However, PGP is proprietary software (its history is quite complex) which therefore limits its use by the open source community and the general public.

  Since the functionality provided by the PGP was useful, eventually the \emph{OpenPGP standard} was developed (\href{https://datatracker.ietf.org/doc/html/rfc4880}{RFC 4880}).
  In July of 2024, \href{https://datatracker.ietf.org/doc/html/rfc9580}{RFC 9580} was published to supersede RFC 4880, but the state of implementations is muddy.
  If you are interested in tech drama, see:
  \begin{itemize}
    \item \href{https://lwn.net/Articles/953797/}{A schism in the OpenPGP world},
    \item \href{https://blog.pgpkeys.eu/critique-critique.html}{A Critique on \enquote{A Critique on the OpenPGP Updates}}.
  \end{itemize}

  A complete and free implementation of the OpenPGP standard is the \emph{GNU Privacy Guard}, abbreviated as GPG (or GnuPG, or GPG2).
  GPG has many uses, but the most \enquote{mainstream}\footnote{This is from my own experience, and is likely not a global truth.} are email encryption and signing\footnote{A nice resource: \url{https://emailselfdefense.fsf.org/en/}}, SSH access (SSH can work with GPG, but should it?), and signing git commits\footnotemark{}.
  \footnotetext{\url{https://docs.github.com/en/authentication/managing-commit-signature-verification/signing-commits}}

  \begin{tcolorbox}
    I believe, for various reasons, e.g. the lack of forward secrecy and allowing users and lack of sensible safeguards\footnote{It is particularly easy to do stupid things with PGP due to the abundance of settings.}, that GPG (and PGP) is legacy tooling and should be deprecated for most purposes.
    While not all of the community would agree with me, I recommend that you avoid PGP unless there are no better alternatives for a specific purpose.
  \end{tcolorbox}

  \begin{tcolorbox}
    Git commits can nowadays also be signed with SSH keys.
    GitHub, GitLab and BitBucket all support it.
    SSH signing was added to Git version 2.34.0 in 2021.
    \href{https://github.com/git/git/blob/master/Documentation/RelNotes/2.34.0.txt}{(changelog)}
  \end{tcolorbox}

  GPG is a command-line tool, however different frontends exist\footnote{\url{https://gnupg.org/software/frontends.html}}.
  The GUI which I have seen used the most is Kleopatra, which is available on Windows, macOS and Linux.
  For Windows, the Kleopatra GUI comes with \href{https://gpg4win.org/index.html}{Gpg4win}.
  Still, I typically use the CLI instead of a GUI (in the rare cases when I have to use GPG).

  For this task, read through the following three articles:
  \begin{itemize}
    \item \href{https://www.digitalocean.com/community/tutorials/how-to-use-gpg-to-encrypt-and-sign-messages}{\textit{How To Use GPG to Encrypt and Sign Messages}} (DigitalOcean)
    \begin{itemize}
      \item Pay special attention to the \enquote{Create a Revocation Certificate} section!
      \item Summarise in a few sentences how public keys are typically shared and made available to other parties.
      
      If you make your key available using this method, can you later make it so that the keys can no longer be searched for / obtained by others?
    \end{itemize}
    \item \href{https://www.gnupg.org/gph/en/manual/x135.html}{\textit{Making and verifying signatures}} (The GNU Privacy Handbook)
    \begin{itemize}
      \item For each of the different signature methods, explain their usefulness and describe a use-case where that specific method is suitable, but the others are not.
    \end{itemize}

    \item \href{https://www.gnupg.org/gph/en/manual/x334.html}{\textit{Validating other keys on your public keyring}}
    \begin{itemize}
      \item Answer the following: what is the difference between \enquote{trust} and \enquote{validity} when speaking about trust in GPG keys and their owners?
    \end{itemize}
  \end{itemize}

  All three articles will help you with the practical tasks.

  \textit{Addendum.}
  If you have to use GPG and want to use it with a YubiKey, check out the article \href{https://support.yubico.com/hc/en-us/articles/360013790259-Using-Your-YubiKey-with-OpenPGP}{\textit{Using Your YubiKey with OpenPGP}}.
\end{task}

\begin{task}{Certificate ToS}
  Read through the short summary on \href{https://www.id.ee/en/article/compete-information-about-the-terms-of-use-of-certificates/\#what-you-need_2}{\textit{What you need to know about the terms and conditions for the use of ID-card certificates}}.

  Then, answer the following questions:
  \begin{itemize}
    \item What must you do when you lose your document or it is stolen?
    \item How can you prevent the unauthorised use of your document?
    \item What is the difference between suspending certificates and declaring them invalid?
    \item What are the differences in revocation and suspension between SK and Zetes ID card certificates?
    \item Can you anonymously use your document digitally?
    \item Does the ID card serve any purpose after its certificates have been invalidated?
  \end{itemize}
\end{task}

\begin{task}{CDOC2}
  Based on \href{https://www.id.ee/en/article/cdoc-2-0/}{\textit{CDOC 2.0 (id.ee)}} answer the following:

  \begin{itemize}
    \item What is the purpose of CDOC?
    \item What are two advantages of the capsule-server approach introduced in CDOC2?
    \item Are capsule servers mandatory in CDOC2?
    \item Is CDOC2 quantum-safe?
    \item Is CDOC2 with capsule servers suitable for long-term encrypted storage of documents?
  \end{itemize}
\end{task}

\begin{task}{ASiC-E containers}
  Task file: \texttt{pubkey.asice}

  Investigate the signature container, then answer the following questions:
  \begin{itemize}
    \item What was the time on my computer at the moment of singing (Estonian local time)?
    \item What is the signing time reported by DigiDoc (Estonian local time)?
    \item Do these times differ? Why yes/why not?
    \item What eID \enquote{device} (ID card, Mobile-ID, Smart-ID) did I use to issue the signature?
    How do you know which?
  \end{itemize}
\end{task}

\newpage
\setcounter{task}{0}

\begin{center}
  Practical tasks
\end{center}

\textbf{Additional instructions}

\begin{itemize}
  \item The GitLab pipeline must succeed before you submit your assignment.
\end{itemize}

\begin{task}{Linking up}
  Task files: \texttt{const1.asice}, \texttt{const2.asice}

  Your task is to create an ASiC-E container with both signatures in it (\autoref{fig:digidoc}).
  Explain what you did to make it so and include commands/screenshots.

  \begin{figure}[h!]
    \center
    \includegraphics[width=\textwidth]{../images/combined}
    \caption{Expected signature status displayed by DigiDoc.}
    \label{fig:digidoc}
  \end{figure}

  Include the file as \texttt{combined.asice} in your repository.

  \textit{Hint.}
  The \href{https://www.id.ee/wp-content/uploads/2021/06/bdoc-spec212-eng.pdf}{BDOC 2.1} reference can help you if you are unable to re-assemble the container in a way such that DigiDoc actually opens it.
\end{task}

\begin{task}{Signing your commits}
  You must set up a commit signing key on TalTech's GitLab.
  I recommend that you use an SSH key for signing commits, but you can also use GPG or S/MIME.

  For the task, you must push a signed commit into your homework branch.
  This commit should include the public key of the key you use for signing: call it simply \texttt{key}.
  The key should not be a binary blob: it should open in a text editor (e.g. PEM).
  Your task is considered complete if the GUI shows that your commit is verified.

  \begin{figure}[h!]
    \center
    \includegraphics[width=\textwidth]{../images/gitlab_verified.png}
    \caption{Verified commit on GitLab.}
  \end{figure}

  Include the full commit link in your report.
\end{task}

\begin{task}{GPG keys}
  Generate a GPG primary key that can only certify other keys, but not sign or encrypt, and which is valid for two months.
  Then, generate a subkey for it which can only sign data, but not encrypt, and which is valid for one month.
  Please do not publish your key to any keyserver online!

  Both keys must use Ed25519.
  For the key metadata, use your real name and your TalTech email.
  For the comment, use \enquote{Assignment 4}.
  \autoref{fig:gpg} displays what your key should look like on the command line.
  You can also find my GPG public key in \texttt{pubkey.asice}.
  Pay attention to the \texttt{[C]} and \texttt{[S]} markers.

  \begin{figure}[h!]
    \center
    \includegraphics[width=\textwidth]{../images/gpg_pubkey}
    \caption{Example metadata for a key with the required parameters.}
    \label{fig:gpg}
  \end{figure}

  Include a screenshot showing that the key is in your keyring in your report (CLI and GUI are both fine).

  Export your full public key in ASCII-armored format, sign it with your Estonian eID\footnote{Smart-ID is fine too, even though it is not officially classified as eID in Estonia.} and include the signed container as \texttt{gpg.asice} in your repository.
  You can leave the container unsigned if you do not have an Estonian eID or do not wish to sign the container for some reason.
  If you leave the container unsigned, I would be interested in the reason, but you do not have to include one.

  Next, generate a revocation certificate for the primary key, if one is not automatically created during key creation.

  \begin{tcolorbox}
    If you ever have to generate a GPG key, make sure to always also generate a revocation certificate!
    Make sure to store this certificate in a safe but available (for you) manner.
    This allows you to revoke a key should you lose it or should it become compromised.
    \tcblower
    Proper key management with GPG is quite complex, e.g. due to the subkey system and the usual complexity of rotating keys.
    As such, revocation certificates become especially important for potential damage control.
  \end{tcolorbox}

  Encrypt the revocation certificate to my ID code using CDOC2 with DigiDoc.
  You can get my ID code from my signed container.
  \begin{tcolorbox}
    In practice, you should of course never leak your revocation certificate, as anyone can permanently revoke your key with it.
  \end{tcolorbox}

  Include the file as \texttt{revocation.cdoc2} in your repository.
\end{task}

\begin{task}{A bit of everything}
  Homework challenge file: \texttt{challenges.txt}

  In the homework challenge file, you will find two columns of data: the first column includes salted hashes of student usernames (e.g. takraa) while the second column includes SHA-256 hashes of registration codes of Estonian organisations.
  Each student is assigned a different organisation.

  Your task is then to:
  \begin{enumerate}
    \item Find which entry corresponds to you (i.e. hashes your student username).
    \item Find the company your challenge is for (i.e. find the SHA-256 pre-image).
    \begin{itemize}
      \item The pre-image is an ASCII string of the form \texttt{itc8280-XXXXXXXX}, where each \texttt{X} represents a decimal digit.
      \item Registration codes in Estonia are 8 decimal digits long.
            The following should help you narrow the search space:
            \begin{itemize}
              \item Registry codes of companies (AS/OÜ) begin with a \enquote{1}.
              \item Registry codes of public institutions begin with a \enquote{7}.
              \item Registry codes of nonprofit organisations (MTÜ) begin with an \enquote{8}.
              \item Registry codes of foundations (sihtasutus) begin with a \enquote{9}.
            \end{itemize}
    \end{itemize}
    \item Find the certificate serial number of any e-Seal certificate of that company.
    \begin{itemize}
      \item Certificates are issued by SK ID Solutions.
      \item I say \enquote{any}, since an organisation might have multiple e-Seal certificates.
      \item NB! Not all certificates are for e-Seals (e.g. some are for encryption).
      \item Do not confuse the certificate serial number with the subject serial (the subject serial is likely the same as the registry code).
      \item The certificate serial has at least 20 decimal digits (in our case).
    \end{itemize}
  \end{enumerate}

  Include in your report:
  \begin{itemize}
    \item The registration code of the company/organisation.
    \item The name of the company.
    \item The certificate serial as a base10 integer.
  \end{itemize}

  Push any code you used to solve the task into GitLab.
  Provide commands and/or screenshots for any steps you performed manually, i.e. that were not solved by your code (if any).
\end{task}

\begin{task}{OCSP}
  Query the OCSP response for the e-Seal certificate from the previous task (use TalTech's if you could not solve it).

  Include the command that you used to query and validate the OCSP response in your report.
  Include the OCSP response (and response only) as \texttt{ocsp.der} in your repository.

  Include also proof (e.g. \texttt{proof.der} or a screenshot in the report) that the intermediate certificate was valid at that time since checking the leaf certificate's OCSP response does not validate the status of the intermediate(s).
\end{task}

\end{document}
